% Figure: the Carnot-engine intermediary used to prove the Clausius inequality.
% A single master reservoir at T0 feeds a reversible Carnot engine that delivers
% heat dQ to the working system at temperature T. Over the system's cycle the
% only net effect is work drawn from a single reservoir, which Kelvin-Planck
% forbids unless the cyclic integral of dQ/T is non-positive.
\begin{figure}[htbp]
\centering
\begin{tikzpicture}[font=\small]
  % master reservoir at top
  \filldraw[engred!18, draw=engred, thick] (0,4.4) rectangle (6,5.4);
  \node[engred] at (3,4.9) {master reservoir at $T_0$};
  % reversible Carnot engine (circle) in the middle
  \draw[engblack, thick, fill=englime!18] (3,3.0) circle (0.7);
  \node[font=\scriptsize, align=center] at (3,3.0) {reversible\\Carnot};
  % working system (rounded rectangle) at bottom
  \filldraw[engblue!14, draw=engblue, thick, rounded corners]
    (1.4,0.4) rectangle (4.6,1.4);
  \node[engblue, align=center] at (3,0.9) {working system\\at $T$ (one cycle)};
  % heat from reservoir into Carnot engine
  \draw[engred, line width=1.4pt, -{Stealth[length=3mm]}] (3,4.4) -- (3,3.7);
  \node[engred, font=\scriptsize, anchor=west] at (3.1,4.05)
    {$\delta Q_0 = T_0\,\delta Q/T$};
  % heat from Carnot engine into working system
  \draw[engorange, line width=1.4pt, -{Stealth[length=3mm]}] (3,2.3) -- (3,1.45);
  \node[engorange, font=\scriptsize, anchor=west] at (3.1,1.9) {$\delta Q$ at $T$};
  % net work delivered out to the side
  \draw[englime!70!engblack, line width=1.4pt, -{Stealth[length=3mm]}]
    (3.7,3.0) -- (5.6,3.0);
  \node[englime!70!engblack, font=\scriptsize, anchor=west] at (5.65,3.0)
    {net work $\delta W$};
  % annotation: Kelvin-Planck bound
  \node[enggray, font=\scriptsize, align=left, anchor=north west] at (-0.3,-0.3)
    {over the system cycle only the reservoir loses heat and work leaves:};
  \node[enggray, font=\scriptsize, anchor=north west] at (-0.3,-0.85)
    {Kelvin-Planck forbids a positive net, so $\oint \delta Q/T \le 0$.};
\end{tikzpicture}
\caption[The Carnot-engine construction behind the Clausius inequality]{The
intermediary that turns the Kelvin-Planck statement into the Clausius inequality.
A single master reservoir at $T_0$ feeds a reversible Carnot engine, which
delivers heat $\delta Q$ to the working system at its local temperature $T$; the
reversible relation fixes the reservoir draw at $\delta Q_0 = T_0\,\delta Q/T$.
Over one full cycle of the working system every state function returns to its
start, so the only net effect of the combined device is that it has drawn
$\oint \delta Q_0 = T_0 \oint \delta Q/T$ from a single reservoir and delivered
an equal net work. The Kelvin-Planck statement forbids net work from a single
reservoir over a cycle, which forces $\oint \delta Q/T \le 0$, with equality for
a reversible cycle. The integrating factor $1/T$ is not guessed; it is forced by
this construction.}
\label{fig:vol04ch01:clausius-construction}
\end{figure}
